\section{Conclusion}
\label{sec:conclusion}

We have presented a method that simultaneously addresses two fundamental failure modes of zero-noise extrapolation: unphysical predictions and unprincipled model selection. By combining constrained optimization (enforcing spectral bounds) with the corrected Akaike Information Criterion (providing bias-corrected model selection), we achieve:

\begin{itemize}
    \item \textbf{Physical guarantee:} 0\% unphysical predictions across all 125 tested configurations, compared to 4\% for standard linear ZNE and up to 12\% at low noise.
    \item \textbf{Competitive accuracy:} MAE of 0.058, matching linear ZNE while outperforming quadratic ZNE by 40\%.
    \item \textbf{Zero-cost automation:} No training data, no hyperparameter tuning, negligible computational overhead ($<$3 seconds).
    \item \textbf{Increasing relevance:} The method becomes more valuable as hardware improves, since lower noise increases the probability of boundary-violating extrapolation.
\end{itemize}

Our key insight is that information-theoretic model selection is \emph{essential} for ZNE in the small-sample regime: without it, bounded optimization alone achieves only 29\% win rate (overfitting within bounds), while AICc selection raises this to 61.6\%. The combination of bounds and AICc is thus synergistic rather than redundant.

\textbf{Future directions.} Several extensions are natural:
\begin{enumerate}
    \item \textbf{Hardware validation} on superconducting processors (Origin Wukong 180, IBM Quantum) to confirm simulator results under realistic noise.
    \item \textbf{Adaptive noise scale factors} selected based on observed decay curvature, reducing the number of required circuit executions.
    \item \textbf{Multi-observable joint estimation} with correlated physical bounds, relevant for variational algorithms estimating Hamiltonians.
    \item \textbf{Integration with dynamical decoupling} and readout mitigation in a unified pipeline, with AutoMitigator selecting the optimal combination.
    \item \textbf{Expanded model family} including Richardson extrapolation and rational functions, with AICc controlling the increased model selection complexity.
    \item \textbf{Scaling studies} on 50--100 qubit circuits to validate performance under spatially correlated noise.
\end{enumerate}

All code, data, and the MCP server are available at \url{https://github.com/mapleleaflatte03/quantum-noise-optimizer} under the MIT license. We encourage integration with existing frameworks (Mitiq, Qiskit) and welcome community contributions.

\begin{acknowledgments}
We thank the Origin Quantum team for providing access to the Wukong 180 calibration data, and the Mitiq and FastMCP open-source communities for their infrastructure. We acknowledge useful discussions with the quantum error mitigation community on the interpretation of physically-bounded extrapolation.
\end{acknowledgments}
